\centerline{\textbf{\large Verification of Ethereum Consensus and Transaction Inclusion}}
\centerline{\textbf{\large with Zero Knowledge Applied to the Transfer of Digital Assets}}

\vspace{1cm}

\rightline{Authors: Javier López - Alejandro Marín}
\vspace{0.3cm}
\rightline{Advisor: PhD Marcos Villagra}

\vspace{0.3cm}

\noindent
\textbf{ABSTRACT}

\vspace{0.3cm}

\noindent
Cross-chain bridges allow assets to be transferred between different blockchains and are a
central element of interoperability. Most of them rely on the lock-and-mint pattern operated by a
trusted relayer, which, after observing a lock of funds on the source chain, mints the equivalent
amount on the destination chain. Since security rests entirely on that operator, bridges have
become one of the most attacked components of the ecosystem. This work designs and implements a
trust-minimized bridge in which the destination contract mints only after verifying
zero-knowledge proofs on-chain. The consensus proof establishes that more than two thirds of the
Ethereum sync committee signed the header from which the execution block is derived, and the
inclusion proof establishes that the lock transaction belongs to that block, calls the lock
contract directly and that its receipt records a successful execution. The bridge does not verify
the finality of the block, and the commitment of each committee is registered by the contract
owner. The relayer ceases to be a trusted minter and becomes a permissionless prover. The
verification logic was implemented in Noir and in Rust for the SP1 virtual machine, with verifiers
in Solidity, and both implementations were compared empirically. Under the measured conditions,
with an earlier version of the Noir consensus circuit that proves a smaller statement, SP1 obtained
the first proof about five times faster, whereas Noir generated each subsequent proof about 5.7
times faster and with less memory. The SP1 time does not include the Groth16 wrapping required for
verification on Ethereum. Under the post-Osaka gas rules, each SP1 verification required about
259,000 gas, between 17 and 18 times less than each of the three Noir proofs. The final Noir version of the
bridge completed real transfers from Ethereum Sepolia to Optimism Sepolia, in about three and a
half minutes each. A lock performed from a smart account through an internal call did not result
in minting because the transaction did not directly target the lock contract and, like its
receipt, exceeded the maximum size supported by the inclusion circuit. A web interface with a comparative mode also shows that the classic
scheme admits a fraudulent mint that the proof-based scheme rejects. The results indicate that
neither tool dominates in every evaluated dimension and that the choice depends on whether the
determining cost is proof generation or on-chain verification.

\vspace{0.3cm}

\noindent
\textbf{Keywords}: cross-chain bridges; zero-knowledge proofs (zk-SNARKs); Noir; UltraHonk; SP1;
on-chain verification; blockchain interoperability; inclusion proof (Merkle-Patricia trie); trust
minimization; Ethereum.
